% Source: https://artofproblemsolving.com/community/c3h1681392_concentric_circles
% Stable identity: aops-topic-1681392-q01
Put the centre at the origin $O$, let the radii be $a<b$, and let the line be $y=d$ with $0<d<a$. Write $\alpha=\sqrt{a^2-d^2}$ and $\beta=\sqrt{b^2-d^2}$, so the intersection points are $(\pm\alpha,d)$ and $(\pm\beta,d)$. The tangent to the circle $x^2+y^2=r^2$ at $(x_0,y_0)$ is $x_0x+y_0y=r^2$.

\begin{asy}
usepackage("amsmath");
size(7cm);
defaultpen(fontsize(10pt));
pair O=(0,0);real a=2,bR=3,d=1.2;
real al=sqrt(a^2-d^2),be=sqrt(bR^2-d^2);
pair A1=(-al,d),A2=(al,d),B1=(-be,d),B2=(be,d);
pair P(pair X,pair Y){return extension(X,X+rotate(90)*X,Y,Y+rotate(90)*Y);}
pair P11=P(A1,B1),P12=P(A1,B2),P21=P(A2,B1),P22=P(A2,B2);
real Rr=a*bR/d;
draw(circle(O,Rr),red+dashed);
draw(circle(O,a),linewidth(0.8));draw(circle(O,bR),linewidth(0.8));
draw((-5.3,d)--(5.3,d),linewidth(0.8));
draw(P11--A1--P12,blue+linewidth(0.5));draw(P21--A2--P22,blue+linewidth(0.5));
draw(P11--B1--P21,deepgreen+linewidth(0.5));draw(P12--B2--P22,deepgreen+linewidth(0.5));
dot(O);dot(A1);dot(A2);dot(B1);dot(B2);dot(P11);dot(P12);dot(P21);dot(P22);
label("$O$",O,S);label("$A_1$",A1,dir(120),fontsize(8pt));label("$A_2$",A2,dir(60),fontsize(8pt));
label("$B_1$",B1,dir(150),fontsize(8pt));label("$B_2$",B2,dir(30),fontsize(8pt));
label("$P_{11}$",P11,SW,fontsize(8pt));label("$P_{22}$",P22,SE,fontsize(8pt));label("$P_{12}$",P12,NE,fontsize(8pt));label("$P_{21}$",P21,NW,fontsize(8pt));
label("\mbox{}",(min(currentpicture,user=true)+max(currentpicture,user=true))/2,black); // keeps the text after the figure black when LaTeX typesets the labels inline
\end{asy}

\textbf{The tangents at $(\alpha,d)$ and $(\beta,d)$.} They are $\alpha x+dy=a^2$ and $\beta x+dy=b^2$. Subtracting, $(\alpha-\beta)x=a^2-b^2=\alpha^2-\beta^2$, so $x=\alpha+\beta$ and $y=\dfrac{a^2-\alpha(\alpha+\beta)}{d}=\dfrac{d^2-\alpha\beta}{d}$. Hence
\[
x^2+y^2=(\alpha+\beta)^2+\dfrac{(d^2-\alpha\beta)^2}{d^2}=\alpha^2+\beta^2+d^2+\dfrac{\alpha^2\beta^2}{d^2}=\dfrac{(d^2+\alpha^2)(d^2+\beta^2)}{d^2}=\dfrac{a^2b^2}{d^2}.
\]
\textbf{The tangents at $(\alpha,d)$ and $(-\beta,d)$.} Replacing $\beta$ by $-\beta$ in this computation gives the intersection $\left(\alpha-\beta,\ \dfrac{d^2+\alpha\beta}{d}\right)$, and the same expansion gives $x^2+y^2=\dfrac{a^2b^2}{d^2}$ again.

The other two points are the mirror images of these two in the $y$-axis (replace $\alpha$ by $-\alpha$). So all four points are at distance $\dfrac{ab}{d}$ from $O$: they lie on the circle with centre $O$ and radius $\boxed{\dfrac{ab}{d}}$, where $d$ is the distance from the centre to the line. $\blacksquare$

(The condition that the line misses the centre is needed: if $d=0$, the tangents at $A_i$ and $B_j$ are parallel.)
